\scanpage{270}
\mspage{11}

\deleted{Soient} \uline{Proposition 4} Soient \deleted{$i<j$} \inserted{$i,j$} dans $I$, \inserted{avec $i<j$.} Conditions équivalentes

a) Il existe une flèche entre $i$ et $j$ (soit $i\to j$, soit $j\to i$)

b) $i$ et $j$ appartiennent à un même intervalle structural \inserted{(\unsure{néc\textsuperscript{t}} unique).}

c) La sous-catégorie pleine de $I$ définie par l'ens d'objets $[i,j]$ est isomorphe à un $\mathbb{\Delta}_{\alpha}$ (on aura $\alpha=j-i$)

On dira alors que $i$ et $j$ sont \uline{directement connectés} dans $I$. (Et on utilisera la $\mathord{=}$ terminologie dans une cat. quelconque : $i$ et $j$ sont dir. conn. \unsure{sss|ssi} $\operatorname{Hom}_I(i,j)$ \uline{ou} $\operatorname{Hom}_I(j,i)$ est non vide.)

\uline{Proposition} 5. Soit $J\subset I$, \deleted{\unsure{consi}} une partie de $I$, considérons la sous-catégorie pleine de $I$ qu'elle définit. Alors

a) Les composantes connexes \inserted{$J_\alpha$} de $J$ sont des intervalles, \deleted{(orientés par l'ordre} et l'ordre induit sur \inserted{un} $J_\alpha$ par celui de $I$ est une orientation (\deleted{\unsure{dite}} \uline{orientation induite}) de $J_\alpha$.

% slanted
\margin{On appelle \uline{sous-intervalle} de $I$, une sous-catégorie pleine connexe, i.e. qui est \unsure{elle-même|elle-$\mathord{=}$} un intervalle. Notons d'ailleurs qu'une sous-catégorie $I'$ de $I$ qui est un intervalle, est nécessairement pleine, cf plus bas \S{} des foncteurs d'intervalles.}

b) \inserted{Supposons \deleted{\unsure{$J\ne\varnothing$}} $\operatorname{card}J\geq1$.} Soit $i$ le plus petit élément de $J$, $j$ le plus grand élément, \deleted{Supposons} donc $i\leq j$. Pour que $J$ lui-même soit un intervalle, il f. et s. que tout sommet critique contenu dans $[i,j]$ soit élément de $J$.

\uline{NB} Si $s$ \inserted{$\in[i,j]$} est un sommet critique \deleted{de $[i,j]$ \unsure{et}} \inserted{de $I$}, \deleted{Si $i,j$ sont deux éléments de $I$ qui} \deleted{qui n'est pas dans $J$, alors} qui n'est pas \deleted{\unsure{dans|en}} $J$, alors par prop. 3 il disconnecterait $J$ \deleted{\unsure{en} $\unsure{2|q}$}, i.e. $J$ serait somme directe de deux ss-catégories $J'$ et $J''$ non vides ($i\in J'$, $j\in J''$), donc $J$ ne serait pas un intervalle.

% archive sheet number: 269
